% エントリ 11(2026-11-07): 集団解析に渡す chi_SR
\section{[E11, 2026-11-07] 集団解析に渡す $\chi_\mathrm{SR}$: 準平衡と自由場をつなぐ}\label{sec:e11}

E10 の準平衡の表(\texttt{chisr.py})に自由場の側をつなぎ、集団解析でそのまま呼べる関数 $\chi_\mathrm{SR}(M,\chi_0,T;\mu,f)$ にした。計算は \texttt{calc/chisr.py} の \texttt{chi\_sr}、照合は \texttt{calc/e11\_chisr.py}、出力は \texttt{calc/out/e11.json}。共同研究者に渡す単体版は \texttt{share/chi\_sr/}(\texttt{calc/export\_chisr.py} が表を作り、プロジェクト版と一致することを確かめる)。

\subsection{つなぎ方}

\paragraph{領域の判定}
BGLS の B と C の境界(式 53、56)にならう。次の二つが、どちらも自由場で飽和する量 $\varepsilon_\mathrm{max}=\chi_0-\chi_c^{(211)}$ より小さければ準平衡、そうでなければ自由場とする。
\begin{itemize}
\item 準平衡の占有数 $\varepsilon_\mathrm{eq}$(E3)
\item 行き過ぎの占有数 $\varepsilon_\mathrm{th}=\sqrt{2\gamma_\mathrm{SR}\ln(GM^2)/\gamma_\mathrm{BH}}$(BGLS の式 58)
\end{itemize}

\paragraph{準平衡の側}
二つの補正を入れて、E10 の表を引く。
\begin{itemize}
\item 1 個の量子から $\varepsilon_\mathrm{eq}$ まで育つ時間 $t_\mathrm{grow}=\ln(\varepsilon_\mathrm{eq}GM^2)/\gamma_\mathrm{SR}$ を、年齢から引く。
\item 行き過ぎで失う $\varepsilon_\mathrm{th}$ を、初期スピンから引く。
\end{itemize}
境界に近づくと $\varepsilon_\mathrm{th}\to\varepsilon_\mathrm{max}$ で、自由場の飽和に連続につながる。

\paragraph{自由場の側}
E2 の半解析と同じ数え方にした。
\begin{itemize}
\item 211 は 1 個の量子から $\ln N_\mathrm{sat}$ 回の $e$-fold で飽和し、臨界スピンまで落ちる。BH の質量の減少を入れる($\alpha\gtrsim0.2$ で臨界スピンが 0.05 ほど下がる)。
\item 322 は $t=0$ から $e$-fold を積む。
\end{itemize}

\paragraph{印}
戻り値に、0(超放射しない)、1(準平衡)、2(自由場)、3(この模型の外)の印をつける。3 は次の二つで、$\chi_0$ をそのまま返す。
\begin{itemize}
\item 準平衡で $\alpha>0.30$。
\item 211 が超放射せず、322 から育つ BH で、$f<10^{16}\GeV$。
\end{itemize}

\subsection{結果}

\begin{figure}[H]
\centering
% E11: 準平衡と自由場をつないだ chi_SR の f 依存(mu = 0.71e-12 eV)。データは calc/e11_chisr.py と e05_kinetics.py
\begin{tikzpicture}
\begin{axis}[note, width=0.72\linewidth, height=0.42\linewidth, xmode=log, xmin=1e10, xmax=1e19, ymin=0.15, ymax=0.82,
  xlabel={$f\ [\mathrm{GeV}]$}, ylabel={年齢 $T$ のあとのスピン},
  legend style={font=\scriptsize, at={(0.02,0.03)}, anchor=south west, fill=white, fill opacity=0.85, text opacity=1, draw=none}, legend cell align=left]
  \fill[cM!12] (axis cs:1e15,0.15) rectangle (axis cs:1e17,0.82);
  \node[font=\scriptsize, color=ink2, anchor=north] at (axis cs:4e15,0.81) {境目(領域 B)};
  \addplot[cA, line width=1.0pt] table[x=f, y=chi_GW241011_T7] {figs/data/e11_scan.dat}; \addlegendentry{GW241011 の主星、$T=10^7$ 年}
  \addplot[cA!60!black, line width=1.0pt, dashed] table[x=f, y=chi_GW241011_T9] {figs/data/e11_scan.dat}; \addlegendentry{GW241011 の主星、$T=10^9$ 年}
  \addplot[cC, line width=1.0pt] table[x=f, y=chi_M15_T9] {figs/data/e11_scan.dat}; \addlegendentry{$15\Msun$、$\chi_0=0.7$、$T=10^9$ 年}
  \addplot[only marks, mark=o, mark size=1.5pt, cB] table[x=f, y=chi_t7] {figs/data/e05_chi_vs_f.dat}; \addlegendentry{E5 の ODE(GW241011、$10^7$ 年)}
\end{axis}
\end{tikzpicture}

\caption{つないだ $\chi_\mathrm{SR}$ の $f$ 依存($\mu=0.71\times10^{-12}\eV$)。丸は E5 の ODE。灰色は準平衡と自由場の境目(BGLS の領域 B)で、自由場の側の 322 の段階が入るかどうかでスピンが飛ぶ。}
\label{fig:e11}
\end{figure}

\begin{result}
\begin{itemize}
\item \textbf{ODE との一致}(GW241011 の主星、$f=10^{11}$--$10^{15}\GeV$、$T=10^6$、$10^7$、$10^8\yr$): 差は最大 $1.6\times10^{-4}$。$15$、$19.6$、$25\Msun$ と $f=10^{12}$、$4\times10^{12}$、$10^{13}\GeV$ では、$10^{10}\yr$ まで小数 3 桁で一致した。
\item \textbf{準平衡の側で 322 の段階は起きない}: ODE は 322 の超放射も含むが、上の例では BH は $10^{10}\yr$ まで 211 の臨界にとどまり、322 の臨界まで落ちない。準平衡では 322 が最初から占有されているので速く育つとも考えられたが、そうはならない。
\item \textbf{自由場の半解析との一致}($f=10^{19}\GeV$、無作為な 400 点): 差の中央値 $3\times10^{-7}$、93.5\% が 0.01 以内、最大 0.034。残りの差は、臨界スピンを解くときのエネルギーの近似(水素原子の主要項か $\alpha^5$ までか)の違いによる。BH の質量の減少を入れる前は最大 0.28 だった。
\item \textbf{境目}: $f\approx10^{16}\GeV$ 付近で、GW241011 の主星の $T=10^9\yr$ のスピンが 0.2 ほど飛ぶ。自由場の側では 322 が飽和するが、準平衡の側では 322 の段階が起きないからである。この領域の運動論は E5 でも解法によって結果が変わったので、確かでない。窓($f\sim10^{12}\GeV$)には効かない。
\item \textbf{速さ}: $10^6$ 個の標本で 0.37 秒。単体版はプロジェクト版と 20 万個の標本で完全に一致した。
\end{itemize}
\end{result}

\begin{caveat}
\begin{itemize}
\item \textbf{重い BH}: 211 が超放射せず 322 から育つ BH の自己相互作用は、模型に入っていない(印 3)。$\mu=0.71\times10^{-12}\eV$ では $\gtrsim35\Msun$ にあたり、GWTC-5 の $45\Msun$ 以上の速く回る群が含まれる。322 と 544 の二準位(\arxiv{2412.03655} の $n\le5$)で同じことをする必要がある。
\item \textbf{自己相互作用の率}: 主要項で、相対論的な補正(E7)は入れていない。
\item \textbf{生まれたときのスピンと年齢}: 分布は集団解析の側の入力である。
\end{itemize}
\end{caveat}
